\begin{lemma}
For every integer \(t\ge2\) there is a constant \(C=C(t)>0\) with the
following property.  Let \(K\) be a finite field of order \(q\), with
\(q\ge C\), and let \(G=G(t,K)\).  If
\[
  k\ge Cq(\log q)^2,
\]
then
\[
  \fwi_k(D^*(G))\le (Cq^t)^k.
\]
\end{lemma}
\begin{proof}
Fix \(t\ge2\).  All constants below are allowed to depend only on \(t\).
Take \(q\) sufficiently large.  By \(\noderef{PolarityGraphParameters}\), the
polarity graph \(G=G(t,K)\) is an \((n,d,\lambda)\)-graph with
\[
  q^t/2\le n\le2q^t,\qquad
  q^{t-1}/2\le d\le2q^{t-1},\qquad
  \lambda\le2\sqrt d. \tag{1}
\]

Let \(\mathcal T\) be the rooted tree whose vertices are forward independent
tuples in \(D^*(G)\), in the sense of \(\noderef{ForwardIndependentTuple}\),
with the empty tuple as root, and whose children are obtained by appending one
more vertex of \(D^*(G)\).  A forward independent tuple
\[
  ((a_1,b_1),\ldots,(a_m,b_m))
\]
is exactly a \(D^*\)-consistent tuple in the sense of
\(\noderef{StarProductConsistentTuple}\): \(a_i b_i\in E(G)\) for every \(i\),
and if \(i<j\) and \(a_i b_j\in E(G)\), then \(a_j b_i\in E(G)\).  This is
just the negation of the arc condition in \(\noderef{StarProductDigraph}\).
When \(G=G(t,K)\), \(\noderef{PolarityGraph}\) identifies the vertices of
\(G\) with projective points and identifies the adjacency relation
\(xy\in E(G)\) with orthogonality \(x\perp y\).

We mark children of each vertex of \(\mathcal T\).  Consider a consistent
tuple
\[
  \sigma=(a_1,b_1,\ldots,a_m,b_m).
\]
For a projective point \(y\), let
\[
  W^\sigma(y)=\operatorname{span}\{b_i:a_i y\in E(G)\},
  \qquad
  r^\sigma(y)=\dim W^\sigma(y).
\]
If a pair \((a,b)\) extends \(\sigma\), then \(ab\in E(G)\), and for every
\(i\) with \(a_i b\in E(G)\), consistency forces \(a b_i\in E(G)\).  Thus
\(a\) is orthogonal to the whole subspace \(W^\sigma(b)\).  Consequently, for
fixed \(b\), the number of possible \(a\)'s is at most
\[
  \frac{q^{t+1-r^\sigma(b)}-1}{q-1}\le 2q^{t-r^\sigma(b)}. \tag{2}
\]
In particular, there are no extensions with \(r^\sigma(b)=t+1\).

For \(0\le r\le t+1\), put
\[
  U_r^\sigma=\{y:r^\sigma(y)\le r\},\qquad
  Z_r^\sigma=U_r^\sigma\setminus U_{r-1}^\sigma,
\]
with \(U_{-1}^\sigma=\emptyset\).  Fix \(0\le r\le t\) with
\(|U_r^\sigma|>0\), and choose \(\ell=\ell_r^\sigma\in\{0,\ldots,r\}\) for
which \(|Z_\ell^\sigma|\) is maximal.  Then
\[
  |Z_\ell^\sigma|\ge |U_r^\sigma|/(r+1)\ge |U_r^\sigma|/(2t). \tag{3}
\]

Call a point \(b\in Z_r^\sigma\) popular if it lies in at least
\(|Z_\ell^\sigma|/(16q)\) subspaces \(W^\sigma(y)\) with
\(y\in Z_\ell^\sigma\).  Since each such \(W^\sigma(y)\) has dimension
\(\ell\), it contains at most \(2q^{\ell-1}\) projective points.  Double
counting pairs \((b,y)\) with \(b\subseteq W^\sigma(y)\) gives at most
\(32q^\ell\) popular points in \(Z_r^\sigma\).  We mark every extension
\((a,b)\) with \(b\) popular.  By (2), this creates at most
\[
  \sum_{r=0}^t 32q^{\ell_r^\sigma}\cdot 2q^{t-r}
  \le 128tq^t \tag{4}
\]
marked children.

Next, for the same \(r\), call a point \(a\) poor if
\[
  |N_G(a)\cap Z_\ell^\sigma|\le |Z_\ell^\sigma|/(8q).
\]
Let \(P_r\) be the set of poor points.  By definition,
\[
  e_G(P_r,Z_\ell^\sigma)\le |P_r||Z_\ell^\sigma|/(8q). \tag{5}
\]
Applying \(\noderef{ExpanderMixingLemma}\) and using the parameter bounds
(1), we also have
\[
  e_G(P_r,Z_\ell^\sigma)
  \ge |P_r||Z_\ell^\sigma|/(4q)
       -4q^{(t-1)/2}\sqrt{|P_r||Z_\ell^\sigma|}. \tag{6}
\]
Combining (5) and (6) yields
\[
  |P_r||Z_\ell^\sigma|\le 1024q^{t+1}. \tag{7}
\]
Since \(|Z_r^\sigma|\le |Z_\ell^\sigma|\), a second application of
\(\noderef{ExpanderMixingLemma}\), again using (1), gives
\[
  e_G(P_r,Z_r^\sigma)\le 5000q^t. \tag{8}
\]
Mark every extension \((a,b)\) for which \(a\) is poor with respect to
\(r^\sigma(b)\).  Summing (8) over \(0\le r\le t\), this contributes at most
\[
  10^4tq^t \tag{9}
\]
additional marked children.

Choose \(A=A(t)\) so large that the sum of the bounds in (4) and (9) is at
most \(Aq^t\).  Thus every vertex of \(\mathcal T\) has at most
\[
  h=Aq^t
\]
marked children.  We now prove that every path has at most
\[
  w=Aq\log q
\]
unmarked vertices, after increasing \(A\) if necessary.

Suppose \(\sigma'=\sigma\cup(a,b)\) is an unmarked child.  Put
\(r=r^\sigma(b)\) and \(\ell=\ell_r^\sigma\).  Since \(a\) is not poor,
\[
  |N_G(a)\cap Z_\ell^\sigma|\ge |Z_\ell^\sigma|/(8q).
\]
Since \(b\) is not popular, it belongs to at most
\(|Z_\ell^\sigma|/(16q)\) of the subspaces \(W^\sigma(y)\) with
\(y\in Z_\ell^\sigma\).  Hence there are at least
\(|Z_\ell^\sigma|/(16q)\) points \(y\in Z_\ell^\sigma\) such that
\(ay\in E(G)\) and \(b\not\subseteq W^\sigma(y)\).  For each such \(y\),
the dimension \(r^\sigma(y)=\ell\) increases to \(\ell+1\) after appending
\((a,b)\).  Therefore
\[
  |U_\ell^{\sigma'}|
  \le |U_\ell^\sigma|-|Z_\ell^\sigma|/(16q). \tag{10}
\]
By (3), \( |U_\ell^\sigma|\le |U_r^\sigma|\le 2t|Z_\ell^\sigma|\).  Thus
\[
  |U_\ell^{\sigma'}|
  \le \left(1-\frac1{32tq}\right)|U_\ell^\sigma|. \tag{11}
\]
Record \(\psi(\sigma')=\ell\) for such an unmarked child.  The same argument
also shows that, whenever \(\psi(\sigma')=\ell\), the preceding set
\(U_\ell^\sigma\) is nonempty.

Now take any root-to-vertex path
\(\sigma_0,\sigma_1,\ldots,\sigma_m\), with \(\sigma_0\) the root.  If it
had more than \(w\) unmarked vertices, then by pigeonhole there would be an
\(\ell\in\{0,\ldots,t\}\) for which at least \(w/(t+1)\) of them satisfy
\(\psi=\ell\).  Along the path the sets \(U_\ell\) are nonincreasing, and at
each such unmarked step they shrink by the factor in (11).  Since
\(|U_\ell^{\sigma_0}|\le n\le2q^t\), choosing \(A(t)\) sufficiently large
gives
\[
  |U_\ell^{\sigma_{i-1}}|
  \le
  2q^t
  \left(1-\frac1{32tq}\right)^{w/(t+1)-1}
  <1
\]
before the last such unmarked step.  Thus \(U_\ell^{\sigma_{i-1}}\) is empty,
contradicting the nonemptiness forced by \(\psi(\sigma_i)=\ell\).  Hence
each path has at most \(w=Aq\log q\) unmarked vertices.

Finally we verify, explicitly, the hypotheses needed to apply
\(\noderef{MarkedTreePathCounting}\).  Let \(P\) be the finite set of
root-to-depth-\(k\) paths in \(\mathcal T\).  Thus an element of \(P\) is a
chain
\[
  p=(\sigma_0,\sigma_1,\ldots,\sigma_k),
\]
where \(\sigma_0\) is the empty tuple and \(\sigma_{j+1}\) is a child of
\(\sigma_j\) for \(0\le j<k\).  Equivalently, \(P\) is the set of ordered
forward independent \(k\)-tuples in \(D^*(G)\): from such a tuple one records
its successive prefixes, and from such a path one reads off the appended
vertices.  This is a bijection by the construction of \(\mathcal T\) and
\(\noderef{ForwardIndependentTuple}\).  Hence
\[
  |P|=\fwi_k(D^*(G))
\]
by \(\noderef{ForwardIndependentTupleCount}\).

The total number of children of any vertex of \(\mathcal T\) is at most
\(|V(D^*(G))|\).  Since \(V(D^*(G))\) is the set of ordered edge-pairs of
\(G\), the parameter bounds in (1) give
\[
  |V(D^*(G))|\le nd\le 4q^{2t-1}
\]
for sufficiently large \(q\).  Put
\[
  \Delta=4q^{2t-1},\qquad h=Aq^t,\qquad w=Aq\log q,
\]
rounding these quantities upward to integers if necessary; increasing \(A\)
and then \(C\) by \(t\)-dependent factors absorbs the ceilings.  Since
\(t\ge2\), choosing \(C\) large enough ensures that, for every \(q\ge C\),
\[
  h=Aq^t\le 4q^{2t-1}=\Delta.
\]
Choosing \(C\ge e\) and \(C\) large compared with \(A\) also gives
\[
  w=Aq\log q\le Cq(\log q)^2\le k,
\]
using the hypothesis \(k\ge Cq(\log q)^2\).  These are the two numerical
hypotheses \(h\le\Delta\) and \(w\le k\) in
\(\noderef{MarkedTreePathCounting}\).

Define the signature map
\[
  \operatorname{sig}:P\to\{0,1\}^k
\]
as follows: for a path \(p=(\sigma_0,\ldots,\sigma_k)\) and an index
\(0\le j<k\), set \(\operatorname{sig}(p)_j=1\) exactly when the child
\(\sigma_{j+1}\) of \(\sigma_j\) is unmarked, and set
\(\operatorname{sig}(p)_j=0\) exactly when it is marked.  By the path bound
proved above, every root-to-depth-\(k\) path has at most \(w\) unmarked
vertices.  In the notation of \(\noderef{BinarySequenceWeight}\), this says
\[
  |\operatorname{sig}(p)|\le w
\]
for every \(p\in P\).

It remains to check the fiber bound for a fixed signature
\(z\in\{0,1\}^k\).  Expose a path \(p\) with \(\operatorname{sig}(p)=z\)
one step at a time.  Suppose a prefix
\((\sigma_0,\ldots,\sigma_j)\) has already been chosen.  If \(z_j=1\), then
the next step is required to be unmarked; there are at most all children of
\(\sigma_j\), and hence at most \(\Delta\), possibilities.  If \(z_j=0\),
then the next step is required to be marked; by the construction of the
marking there are at most \(h\) possibilities.  Therefore, after all \(k\)
steps have been exposed, the number of paths with signature \(z\) is at most
\[
  \Delta^{|z|}h^{k-|z|},
\]
where \(|z|\) is the number of \(1\)'s in \(z\), again in the sense of
\(\noderef{BinarySequenceWeight}\).  This is precisely the per-signature
fiber hypothesis in \(\noderef{MarkedTreePathCounting}\).

Applying \(\noderef{MarkedTreePathCounting}\) to \(P\) with this signature map
and the above values of \(\Delta,h,w\) gives
\[
  \fwi_k(D^*(G))=|P|
  \le
  2^k(4q^{2t-1})^w(Aq^t)^{k-w}.
\]
For \(q\ge C\), after increasing \(C\) so that \(q\ge4\) and \(Aq^t\ge1\),
the right hand side is at most
\[
  2^k q^{2tw}(Aq^t)^k.
\]
Finally, because \(w=Aq\log q\), another \(t\)-dependent increase of \(C\)
ensures
\[
  q^{2tw}\le2^k
\]
whenever \(k\ge Cq(\log q)^2\).  Hence
\[
  \fwi_k(D^*(G))\le4^k(Aq^t)^k\le(Cq^t)^k,
\]
after increasing \(C\) once more so that \(C\ge4A\).  This proves the lemma.
\end{proof}
